% ===========================================================================
% Category: Large Language Model
% Generated from ML_Lab_Manual_Completed.md - edit the Markdown and re-run
% scripts/build_latex_manual.py instead of editing this file by hand.
% ===========================================================================

\chapter{Experiment 18: Large Language Model (LLM)
Experiment}\label{experiment-18-large-language-model-llm-experiment}

\textbf{Category:} Generative / foundation models

\section{1. Objective}\label{objective}

Demonstrate a reproducible LLM workflow: tokenization, inference,
controlled generation, and fine-tuning on a real public classification
dataset.

\section{2. Dataset Used}\label{dataset-used}

DistilBERT (SST-2) for tokenization/sentiment, DistilGPT2 for
generation, and the SMS Spam Collection
(\passthrough{\lstinline!data/sms\_spam.csv!}, 5,572 labeled messages)
for fine-tuning. A balanced subset (300 ham + 300 spam) is split 400
train / 200 stratified test.

\section{3. Required Libraries}\label{required-libraries}

transformers, torch (CUDA), pandas

\section{4. Brief Theory}\label{brief-theory}

LLMs process text as subword tokens (BPE/WordPiece) with an attention
mask; self-attention computes Attention(Q,K,V) = softmax(QKᵀ/√d\_k +
mask)V. Causal generation predicts one token at a time; temperature
scales logits and top-p restricts sampling to the nucleus. Fine-tuning
reuses pretrained weights and replaces the classification head.

\section{5. Procedure}\label{procedure}

\begin{enumerate}
\def\labelenumi{\arabic{enumi}.}
\tightlist
\item
  Tokenize a sentence; inspect tokens, IDs and attention mask.
\item
  Run the sentiment pipeline on fixed test sentences.
\item
  Generate text with greedy decoding and with temperature/top-p
  sampling.
\item
  Load SMS Spam; build the balanced subset and stratified split.
\item
  Fine-tune DistilBERT (fresh 2-class head, AdamW lr 5e-5, 2 epochs,
  batch 16).
\item
  Evaluate accuracy/F1 on the untouched test messages and inspect
  examples.
\end{enumerate}

\section{6. Reference Implementation}\label{reference-implementation}

\begin{lstlisting}[language=Python]
ft_model = AutoModelForSequenceClassification.from_pretrained(model_id, num_labels=2, ignore_mismatched_sizes=True)
optimizer = torch.optim.AdamW(ft_model.parameters(), lr=5e-5)
outputs = ft_model(input_ids=ids, attention_mask=mask, labels=labels)
outputs.loss.backward(); optimizer.step()
\end{lstlisting}

\section{7. Results}\label{results}

{\def\LTcaptype{none} % do not increment counter
\begin{longtable}[]{@{}
  >{\raggedright\arraybackslash}p{(\linewidth - 2\tabcolsep) * \real{0.5000}}
  >{\raggedright\arraybackslash}p{(\linewidth - 2\tabcolsep) * \real{0.5000}}@{}}
\toprule\noalign{}
\begin{minipage}[b]{\linewidth}\raggedright
Item
\end{minipage} & \begin{minipage}[b]{\linewidth}\raggedright
Value
\end{minipage} \\
\midrule\noalign{}
\endhead
\bottomrule\noalign{}
\endlastfoot
Fine-tune loss & 0.705 → 0.079 (2 epochs) \\
Held-out accuracy & 0.9600 \\
Held-out F1 (spam) & 0.9592 \\
Examples & spam 99.4\%, ham 99.5\%, spam 89.7\% confidence (all
correct) \\
\end{longtable}
}

Generation: greedy produced a safe deterministic continuation; nucleus
sampling (T=0.7, p=0.9) produced a different but coherent sentence.
Sentiment scores were in the 90-99\% range on clear sentences.

\section{8. Discussion and Conclusion}\label{discussion-and-conclusion}

The experiment shows the complete practical LLM pipeline and that
fine-tuning a small pretrained model on 400 real messages already gives
96\% spam detection. Limitations: the balanced subset ignores the
natural 13\% spam prior; the models are small and can hallucinate;
fine-tuning data and prompts must be reproducible. No private data is
used.

\section{9. Answers to Viva Questions}\label{answers-to-viva-questions}

\begin{itemize}
\tightlist
\item
  \textbf{What is tokenization?} Splitting text into subword units that
  map to integer vocabulary IDs, letting a fixed vocabulary represent
  any word; special tokens and attention masks handle sequence
  framing/padding.
\item
  \textbf{Pretraining vs fine-tuning vs prompting?} Pretraining learns
  general language representation from massive corpora; fine-tuning
  adapts those weights to a task with labeled data; prompting steers a
  frozen model through instructions/examples.
\item
  \textbf{Why is LLM output not automatically ground truth?} The model
  generates the most plausible continuation, not verified facts; it can
  hallucinate, and confidence scores are not calibrated correctness
  probabilities.
\item
  \textbf{How can prompt wording affect results?} Meaning, tone, format
  and label wording all change the conditional distribution the model
  samples from; small prompt changes can flip outputs, so prompts must
  be fixed and reported.
\end{itemize}

\begin{center}\rule{0.5\linewidth}{0.5pt}\end{center}

\section{10. Complete Code Listing}\label{complete-code-listing}

\emph{Full runnable program for this experiment, extracted from
\passthrough{\lstinline!notebooks/10\_large\_language\_model.ipynb!}
(executed; results above are its actual output).}

\textbf{Listing 18.1 --- Core numerical and plotting libraries}

\begin{lstlisting}[language=Python]
# ---- Core numerical and plotting libraries ----
import os
import numpy as np
import pandas as pd
import matplotlib.pyplot as plt
import seaborn as sns

# ---- PyTorch: used for the fine-tuning demo ----
import torch
import torch.nn as nn
from torch.utils.data import DataLoader, Dataset
# ---- Dataset splitting + metrics for the fine-tune evaluation ----
from sklearn.model_selection import train_test_split
from sklearn.metrics import accuracy_score, f1_score
# ---- Hugging Face Transformers: tokenizers, models, pipelines, generation configs ----
from transformers import (
    AutoTokenizer,
    AutoModelForSequenceClassification,
    AutoModelForCausalLM,
    GenerationConfig,
    pipeline
)

# Styling setup (consistent look across all notebooks)
plt.style.use('seaborn-v0_8-whitegrid' if 'seaborn-v0_8-whitegrid' in plt.style.available else 'default')
plt.rcParams['font.sans-serif'] = 'DejaVu Sans'
plt.rcParams['figure.figsize'] = (10, 5)
plt.rcParams['figure.dpi'] = 110

torch.manual_seed(42)
torch.cuda.manual_seed_all(42)  # reproducible GPU training
device = torch.device('cuda' if torch.cuda.is_available() else 'cpu')  # GPU if available
print(f"PyTorch on {device}")
\end{lstlisting}

\textbf{Listing 18.2 --- Load a pre-trained DistilBERT tokenizer (fast,
subword/WordPiece based)}

\begin{lstlisting}[language=Python]
# ---- Load a pre-trained DistilBERT tokenizer (fast, subword/WordPiece based) ----
model_id = "distilbert/distilbert-base-uncased-finetuned-sst-2-english"
print(f"Loading Tokenizer: {model_id}...")
tokenizer = AutoTokenizer.from_pretrained(model_id)

sample_text = "Machine learning and deep neural networks are transforming modern science!"
# Tokenize with padding/truncation so the output has a fixed shape (here 20 tokens)
encoded = tokenizer(
    sample_text,
    padding="max_length",
    max_length=20,
    truncation=True,
    return_tensors="pt"
)

# Convert integer IDs back to readable subword tokens
tokens = tokenizer.convert_ids_to_tokens(encoded['input_ids'][0])
print(f"Original Text: {sample_text}\n")
print(f"Subword Tokens ({len(tokens)}):\n", tokens)
print(f"Input IDs:      \n", encoded['input_ids'][0].tolist())
print(f"Attention Mask: \n", encoded['attention_mask'][0].tolist())
\end{lstlisting}

\textbf{Listing 18.3 --- Load a pre-trained sentiment-analysis pipeline
(DistilBERT SST-2)}

\begin{lstlisting}[language=Python]
# ---- Load a pre-trained sentiment-analysis pipeline (DistilBERT SST-2) ----
classifier = pipeline("sentiment-analysis", model=model_id, device=0 if torch.cuda.is_available() else -1)

test_sentences = [
    "The new neural network model demonstrated extraordinary accuracy and remarkable speed.",
    "The dataset was poorly labeled, noisy, and the model completely failed to converge.",
    "The gradient descent algorithm performed reasonably well, meeting standard expectations.",
    "A catastrophic failure occurred during training due to exploding gradient issues."
]

# One forward pass per sentence -> label + calibrated confidence score
predictions = classifier(test_sentences)

# Collect results in a table for display
results_table = []
for sent, pred in zip(test_sentences, predictions):
    results_table.append({
        "Sentence": sent,
        "Predicted Label": pred['label'],
        "Confidence Score": f"{pred['score']*100:.2f}%"
    })

df_preds = pd.DataFrame(results_table)
display(df_preds)
\end{lstlisting}

\textbf{Listing 18.4 --- Load a small causal LM (DistilGPT2) for text
generation}

\begin{lstlisting}[language=Python]
# ---- Load a small causal LM (DistilGPT2) for text generation ----
gen_model_id = "distilbert/distilgpt2"
print(f"Loading Text Generator: {gen_model_id}...")
# clean_up_tokenization_spaces=False keeps BPE spacing intact
generator = pipeline("text-generation", model=gen_model_id,
                     device=0 if torch.cuda.is_available() else -1, clean_up_tokenization_spaces=False)

prompt = "Artificial Intelligence will fundamentally transform"

# 1. Greedy Search: always pick the most likely next token (deterministic)
greedy_output = generator(
    prompt, generation_config=GenerationConfig(max_new_tokens=40, do_sample=False)
)[0]['generated_text']

# 2. Temperature + Top-p (Nucleus) Sampling: sample from the most likely tokens up to cumulative p
sampled_output = generator(
    prompt, generation_config=GenerationConfig(max_new_tokens=40, do_sample=True, temperature=0.7, top_p=0.9)
)[0]['generated_text']

print(f"PROMPT: '{prompt}'\n")
print(f"--- 1. GREEDY SEARCH OUTPUT ---\n{greedy_output.strip()}\n")
print(f"--- 2. NUCLEUS SAMPLING (T=0.7, p=0.9) OUTPUT ---\n{sampled_output.strip()}")
\end{lstlisting}

\textbf{Listing 18.5 --- Real dataset: SMS Spam Collection (5,572
labeled messages)}

\begin{lstlisting}[language=Python]
# ---- Real dataset: SMS Spam Collection (5,572 labeled messages) ----
sms = pd.read_csv('../data/sms_spam.csv')
sms['y'] = (sms['label'] == 'spam').astype(int)
print(f"SMS Spam Collection: {len(sms)} messages | ham {(sms.y == 0).sum()}, spam {(sms.y == 1).sum()}")

# Balanced subset for a fast fine-tune demo: 300 ham + 300 spam
ham = sms[sms.y == 0].sample(n=300, random_state=42)
spam = sms[sms.y == 1].sample(n=300, random_state=42)
subset = pd.concat([ham, spam]).sample(frac=1, random_state=42).reset_index(drop=True)

# Stratified split keeps the ham/spam ratio identical in train and test
train_df, test_df = train_test_split(subset, test_size=200, random_state=42, stratify=subset.y)
print(f"Train messages: {len(train_df)} | Test messages: {len(test_df)} | spam share: {subset.y.mean():.1%}")

class SMSDataset(Dataset):
    """Tokenize SMS texts and expose them as PyTorch samples."""
    def __init__(self, texts, labels, tokenizer):
        self.encodings = tokenizer(list(texts), padding=True, truncation=True, max_length=64, return_tensors="pt")
        self.labels = list(labels)

    def __getitem__(self, idx):
        item = {key: val[idx] for key, val in self.encodings.items()}
        item['labels'] = torch.tensor(self.labels[idx], dtype=torch.long)
        return item

    def __len__(self):
        return len(self.labels)

train_loader = DataLoader(SMSDataset(train_df.text, train_df.y, tokenizer), batch_size=16, shuffle=True)
test_loader = DataLoader(SMSDataset(test_df.text, test_df.y, tokenizer), batch_size=32)

# ---- Load the pre-trained classifier with a fresh 2-class head ----
# ignore_mismatched_sizes allows replacing the original SST-2 head with a new one.
ft_model = AutoModelForSequenceClassification.from_pretrained(model_id, num_labels=2, ignore_mismatched_sizes=True)
ft_model.to(device)
optimizer = torch.optim.AdamW(ft_model.parameters(), lr=5e-5)  # small LR for fine-tuning

print("Pre-trained Transformer loaded for fine-tuning!")
\end{lstlisting}

\textbf{Listing 18.6 --- Fine-tuning loop (2 epochs on 400 real SMS
messages)}

\begin{lstlisting}[language=Python]
# ---- Fine-tuning loop (2 epochs on 400 real SMS messages) ----
epochs = 2
ft_model.train()
epoch_losses = []

for epoch in range(epochs):
    running_loss = 0.0
    for batch in train_loader:
        optimizer.zero_grad()

        # Move the batch to the selected device
        input_ids = batch['input_ids'].to(device)
        attention_mask = batch['attention_mask'].to(device)
        labels = batch['labels'].to(device)

        # Forward pass: the model computes the classification loss internally
        outputs = ft_model(input_ids=input_ids, attention_mask=attention_mask, labels=labels)
        loss = outputs.loss
        loss.backward()      # backpropagate through the whole transformer
        optimizer.step()     # update weights

        running_loss += loss.item() * len(input_ids)   # de-average the batch loss

    avg_loss = running_loss / len(train_df)
    epoch_losses.append(avg_loss)
    print(f"Epoch [{epoch+1}/{epochs}] Fine-Tuning Loss: {avg_loss:.4f}")
\end{lstlisting}

\textbf{Listing 18.7 --- Evaluate the fine-tuned model on held-out real
SMS messages}

\begin{lstlisting}[language=Python]
# ---- Evaluate the fine-tuned model on held-out real SMS messages ----
ft_model.eval()
preds, targets = [], []

with torch.no_grad():
    for batch in test_loader:
        # Move the batch to the same device as the fine-tuned model (GPU when available)
        input_ids = batch['input_ids'].to(device)
        attention_mask = batch['attention_mask'].to(device)
        logits = ft_model(input_ids=input_ids, attention_mask=attention_mask).logits
        preds += logits.argmax(-1).tolist()       # predicted class (0 = ham, 1 = spam)
        targets += batch['labels'].tolist()

print(f"Held-out test accuracy: {accuracy_score(targets, preds):.4f}")
print(f"Held-out test F1 (spam): {f1_score(targets, preds):.4f}")

# ---- Show a few individual predictions with confidence ----
labels_map = {0: "ham", 1: "spam"}
with torch.no_grad():
    for i in range(3):
        message = test_df.text.iloc[i]
        enc = tokenizer(message, truncation=True, max_length=64, return_tensors="pt").to(device)
        probs = torch.softmax(ft_model(**enc).logits, dim=1).cpu().numpy()[0]
        print(f"\nMessage: {message[:75]}...")
        print(f"True: {labels_map[test_df.y.iloc[i]]} | Predicted: {labels_map[probs.argmax()]} ({probs.max()*100:.1f}%)")
\end{lstlisting}
